\section{InstructGPT y el primer ChatGPT: efectos, costos y límites}

Para saber si RLHF vale la pena no basta con mirar la función de pérdida: también hay que ver qué prefieren los usuarios, cuántos recursos se invierten y si las capacidades previas retroceden. Esta sección examina, en orden, tres evidencias presentadas en clase: la curva de preferencias, el costo de entrenamiento y el balance de la experiencia con ChatGPT a principios de 2023.

\subsection{100 veces más pequeño, ¿por qué sigue siendo preferido?}

El artículo de InstructGPT somete modelos base GPT-3, SFT, PPO y PPO-ptx de distintos tamaños a la misma evaluación de preferencias humanas. El eje horizontal es el tamaño del modelo en parámetros y el vertical, la frecuencia con que se elige frente a cierta línea base. El resultado más llamativo es que, en esa distribución de customer-prompt, la variante de InstructGPT de 1.3B puede preferirse sobre el GPT-3 base de 175B.

\lecturefigure{slide-08-instructgpt-preferences.jpg}{Curva de preferencias de InstructGPT: el alignment fine-tuning puede cambiar de forma notable la utilidad que perciben los usuarios.}{Video oficial de Stanford Online 00:09:24--00:15:18.}

\begin{knowledgebox}{Lectura de la figura: primero el eje y, después el tamaño del modelo}
El eje y no mide intelligence general ni el promedio de todos los benchmarks, sino la preference win rate humana sobre una distribución específica de prompts. La figura respalda que «un aligned model más pequeño puede ajustarse mejor a la intención del usuario»; no respalda que «el modelo de 1.3B supera al de 175B en conocimiento, razonamiento y todas las tareas». La curva de PPO-ptx recuerda además que la ganancia en preferencia y la conservación de capacidades deben medirse al mismo tiempo.
\end{knowledgebox}

Si el experimento de preferencias se formula como una variable aleatoria binomial, la tasa de victorias estimada del modelo $A$ frente al modelo $B$ es

\[
\widehat{p}_{A>B}=\frac{1}{N}\sum_{i=1}^{N}\mathbb{1}[A\text{ se elige en la comparación }i],
\]

donde $N$ es el número de comparaciones válidas y $\mathbb{1}[\cdot]$ es la función indicadora. Este estadístico depende de la distribución de prompts, de los labelers, del orden de presentación y de los criterios de evaluación; con otro conjunto de tareas u otra población, $\widehat{p}_{A>B}$ puede cambiar.

\begin{warningbox}{Preference no es sinónimo de capability}
La preferencia puede premiar la claridad, la cooperación, el rechazo de solicitudes indebidas y el apego al formato, pero también la seguridad aparente, la complacencia o un estilo familiar. Por eso, la preference win debe considerarse junto con la veracidad, la calibración, la robustez, la exactitud en tareas especializadas y las evaluaciones de seguridad, y no puede cargar por sí sola con todo el significado de «el modelo es mejor».
\end{warningbox}

\subsection{El balance de recursos del alignment fine-tuning}

La slide de costos coloca en un mismo eje vertical el compute de GPT-3 pretraining, de InstructGPT SFT y de InstructGPT RL, y muestra aparte, a la derecha, la retroalimentación humana. Lo que expone no es que «el alignment sea gratis», sino que las distintas dimensiones de recursos son asimétricas: frente al preentrenamiento, el cómputo de SFT/RL es muy pequeño; pero reunir demonstrations y comparisons de alta calidad exige mucho tiempo humano, diseño de criterios y control de calidad.

\lecturefigure{slide-09-instructgpt-training-costs.jpg}{Costo de entrenamiento de InstructGPT: el alignment compute es muy pequeño, pero la retroalimentación humana aún requiere unas 20,000 horas.}{Video oficial de Stanford Online 00:15:18--00:19:05.}

\begin{knowledgebox}{Lectura de la figura: no sumes directamente PFLOP/s-days y dólares}
La gráfica de barras de la izquierda compara cómputo; el texto de la derecha da la estimación del ponente para la retroalimentación humana: unos 500,000 dólares y unas 20,000 horas. Son recursos distintos: el compute puede paralelizarse con hardware, mientras que el juicio humano está sujeto a restricciones de reclutamiento, capacitación, consistencia y representatividad. El costo real del sistema incluye también la gobernanza de datos, el diseño de evaluaciones, la iteración de ingeniería y los experimentos fallidos; la slide no afirma cubrir el costo de todo el ciclo de vida.
\end{knowledgebox}

\begin{table}[H]
\centering
\caption{Tres tipos de recursos que no pueden reducirse a una sola cifra de precio}
\begin{tabularx}{\textwidth}{L{2.9cm} X X}
\toprule
Recurso & Función principal & Riesgo principal \\
\midrule
Pretraining compute & Adquirir capacidades amplias de lenguaje y de tareas & Costo alto, límites de datos opacos, fecha de corte del conocimiento \\
Alignment compute & Hacer las capacidades existentes más controlables y más acordes con los objetivos de interacción & Reward hacking, retroceso de capacidades, sobreajuste a las preferencias \\
Human feedback & Aportar demonstrations, comparisons y juicios sobre errores & Representatividad insuficiente, deriva de criterios, fatiga y problemas de consistencia \\
\bottomrule
\end{tabularx}
\end{table}

\subsection{ChatGPT: la interfaz de diálogo mejora la usabilidad, pero no elimina el desajuste}

El ponente presenta ChatGPT como el siguiente paso de InstructGPT: el diálogo se vuelve una interfaz general, el usuario puede repreguntar, corregir y pedir que se reescriba una respuesta; también mejora el rechazo de tareas dañinas. Pero eso no significa que «la alineación esté terminada». A principios de 2023, las limitaciones más notorias seguían siendo la hallucination y la prompt sensitivity: el modelo inventa hechos, y distintas formulaciones pueden provocar respuestas de distinta calidad o comportamientos de seguridad distintos.

\lecturefigure{slide-10-chatgpt-lessons.jpg}{Balance de ChatGPT presentado en clase a principios de 2023: dialog, tareas complejas, hallucination y prompt sensitivity.}{Video oficial de Stanford Online 00:19:05--00:20:19.}

\begin{warningbox}{Límite temporal: no expliques esta slide con el ChatGPT posterior}
Esta página describe el sistema público del 2023-01-17. En la sesión de Q\&A, el ponente dice explícitamente que la cantidad exacta de datos y los pasos de entrenamiento de ChatGPT no se han publicado, y que solo puede afirmarse que el proceso es, en general, similar al de InstructGPT; las cifras que recuerda de comparisons, demonstrations y las 20,000 horas también son aproximadas. Los apuntes no deben atribuir retroactivamente a ese momento, como especificaciones conocidas, las capacidades posteriores del producto, el uso de herramientas ni los cambios de política.
\end{warningbox}

\begin{importantbox}{Por qué la prompt sensitivity también es un problema de alignment}
Si dos prompts semánticamente equivalentes provocan, solo por diferencias de redacción, niveles completamente distintos de veracidad, rechazo o calidad de razonamiento, lo que el modelo aprendió sigue siendo un patrón superficial y no una intención de tarea estable. Un sistema ideal debería conservar una personalización razonable y, a la vez, mantener sin cambios su comportamiento esencial ante transformaciones que no alteran el significado, además de preguntar activamente cuando haya ambigüedad.
\end{importantbox}

\subsection{Resumen de la sección}

InstructGPT demuestra que el alignment fine-tuning puede mejorar de forma notable la preferencia humana en una distribución específica, con un costo de cómputo muy inferior al del preentrenamiento; pero el trabajo de retroalimentación, la representatividad y el riesgo de proxy siguen siendo costosos. El primer ChatGPT demostró además que la interfaz de diálogo por sí misma mejora la usabilidad, y a la vez dejó ver que la hallucination y la prompt sensitivity siguen sin resolverse.
